\chapter{Data Description}

The datasets used in this project consist of historical financial time series data from two major stock market indices, namely the NIFTY50 index of India and the SSE50 index of China. These indices are widely regarded as important indicators of the economic and financial performance of their respective markets. The NIFTY50 index represents the weighted performance of fifty leading companies listed on the National Stock Exchange of India, while the SSE50 index tracks fifty large scale and highly liquid companies listed on the Shanghai Stock Exchange. Since both indices reflect the behaviour of major financial markets, they provide suitable benchmark datasets for evaluating stochastic forecasting models.

The historical data were collected over multiple years with observations recorded on trading days. Each observation in the datasets originally contained multiple financial attributes including opening price, highest price, lowest price, closing price, trading volume and percentage change for the corresponding trading day. These variables collectively provide different perspectives regarding daily market activity, volatility and trading behaviour. The opening price represents the initial market value at the start of a trading session, while the highest and lowest prices indicate the range of price movement during the day. Trading volume reflects the number of transactions executed, and percentage change measures the relative movement of prices across consecutive trading sessions.

Although the datasets contained several financial indicators, the primary focus of this project was on modeling and forecasting the daily closing prices of the indices. The closing price is commonly regarded as one of the most informative indicators in financial analysis because it summarizes the final market consensus at the end of each trading session. Consequently, the remaining variables such as opening price, highest price, lowest price, volume and percentage change were excluded from the final modeling process. This reduction in dimensionality was performed in order to concentrate specifically on the temporal dynamics of the closing price series and to simplify the learning framework for the stochastic neural differential equation models considered in this study.

Financial time series data generally exhibit non-stationary behaviour characterized by long term trends, volatility clustering, irregular fluctuations and sudden transitions. Such characteristics make financial forecasting particularly challenging because the statistical properties of the data may change over time. Moreover, stock market movements are strongly influenced by external economic conditions, government policies, geopolitical events and investor sentiment, all of which contribute to uncertainty and stochastic behaviour within the series. Due to these complex properties, financial datasets provide an appropriate framework for studying stochastic neural differential equation models which are specifically designed to capture uncertainty and nonlinear temporal dynamics.

Before constructing the forecasting models, the datasets were organized and processed in chronological order to preserve the natural temporal structure of the financial series. Since time series forecasting requires learning from past observations to predict future behaviour, maintaining sequential consistency is essential for realistic modeling. Missing observations and non trading days were handled implicitly through the ordered market records available in the datasets. The resulting financial series therefore represented continuous sequences of historical market behaviour suitable for supervised learning applications.

For the forecasting task, the data were transformed into sequential input-output pairs using sliding window based approaches. In this method, a fixed number of previous observations were utilized as input sequences for predicting future index values. The sliding window framework enables the neural networks to learn temporal dependencies and evolving patterns from historical observations. By continuously shifting the observation window across the dataset, multiple training samples were generated from the same financial series, thereby improving the learning capability of the forecasting models.

In this project, prediction horizons of 1, 3, 7 and 14 days were considered in order to evaluate both short term and relatively longer term forecasting performance. The use of multiple forecasting horizons allows analysis of how predictive accuracy changes as the forecasting interval increases. Short term forecasting generally focuses on capturing immediate market behaviour and local temporal dependencies, whereas longer forecasting horizons involve greater uncertainty and increased difficulty due to the accumulation of stochastic effects over time. Evaluating the models under different horizons therefore provides a more comprehensive understanding of their forecasting capability and robustness.

Normalization techniques were applied to the financial series before feeding them into the neural network models in order to improve numerical stability and facilitate efficient optimization during training. Since raw financial values may vary substantially across different periods and market conditions, normalization helps prevent instability in gradient based learning algorithms and allows the models to converge more effectively. The datasets were subsequently divided into training and testing subsets based on chronological order, ensuring that the temporal dependency structure of the series was preserved. Such a division reflects realistic forecasting conditions in financial applications, where future information remains unavailable during the training phase.

The use of financial datasets from two different markets also allows comparative analysis of the robustness and adaptability of stochastic differential equation based neural networks under varying economic and market conditions. The Indian and Chinese financial markets possess different structural and behavioural characteristics, which may influence the stochastic dynamics present within the data. Since both indices contain complex nonlinear and uncertain patterns, they provide an effective framework for studying the capability of Brownian motion driven SDE-Net models and Lévy driven neural differential equation models in capturing uncertainty and forecasting future market behaviour.